\input{../../../stock/en-tete_v6.tex}

\newcommand{\titrechapitre}{TD d'algorithmique}
\newcommand{\numerochapitre}{1}

\renewcommand{\headrulewidth}{0.4pt}
\renewcommand{\footrulewidth}{0.4pt}
\fancyfoot[L]{Notes de Raphaël JONTEF}
\fancyfoot[C]{\thepage}
\fancyfoot[R]{Enseirb-Matmeca}
\fancyhead[L]{Cours de Mathématiques}
\fancyhead[R]{\titrechapitre}

\begin{document}
\input{../../../stock/commands.tex}

% Page de titre custom
\setlength{\headheight}{15pt}
\begin{titlepage}
    \centering
    \vspace*{3cm}
    {\huge Chapitre \numerochapitre\par}
    {\Huge \bfseries\titrechapitre\par}
    \vspace{2cm}
    {\Large Notes rédigées par Raphaël JONTEF\par}
    \vspace{0.5cm}
    {\normalsize Cours enseigné par François DUFOUR\par}
    \vspace{4cm}
    {\small Enseirb-Matmeca\par}
    {\small filière informatique\par}
    \vfill
    {\large \today\par}
\end{titlepage}

% plan du cours
\tableofcontents
\newpage


% -----------Corps du document--------------

\section*{Séance 6}
\begin{proposition}{première}{un équivalent}
    On a~:
    $$\frac{(2n)!}{(n!)^2} \symbolelim{n \to + \infty}{}{\sim} \frac{4^n}{\sqrt{\pi n}}$$
\end{proposition}

\begin{demonstration}
    On applique la formule de Stirling~:
    \begin{align*}
        (2n)! &\symbolelim{n \to + \infty}{}{\sim} \sqrt{4 \pi n} \left(\frac{2n}{e}\right)^{2n} \\
        (n!)^2 &\symbolelim{n \to + \infty}{}{\sim} 2 \pi n \left(\frac{n}{e}\right)^{2n}
    \end{align*}
    On en déduit~:
    $$\frac{(2n)!}{(n!)^2} \symbolelim{n \to + \infty}{}{\sim} \frac{4^n}{\sqrt{\pi n}}$$
\end{demonstration}

\begin{proposition}{finale}{équivalent de la quantité en jeu}
    On a~:
    $$\binom{2n-1}{n-1} \symbolelim{n \to + \infty}{}{\sim} \frac{1}{2\sqrt{\pi}} \frac{4^n}{\sqrt{n}}$$
\end{proposition}

\begin{demonstration}
    On a~:
    \begin{align*}
        \binom{2n-1}{n-1} &= \frac{(2n-1)!}{n!(n-1)!}\\
        &= \frac{n}{2n} \frac{(2n)!}{(n!)^2} \\
        &\symbolelim{n \to + \infty}{}{\sim} \frac{1}{2} \frac{4^n}{\sqrt{\pi n}} \\
    \end{align*}
    D'où le résultat et le caractère exponentiel en $n$ de cette quantité.
\end{demonstration}

\section*{Séance 7}

\subsection*{Exercice 1}
\subsubsection*{Question 1}
Considérons par l'absurde un algorithme $\mc{A}$ de complexité en temps $C_n$ strictement inférieure à un $\mc{O}(n)$, tel que pour toute entrée \code{T} de taille $n$, $\mc{A}(\code{T})$ lit strictement plus de 1\% du tableau.\\
Par définition de la stricte domination~:
$$\frac{C_n}{n} \flechelim{n \to + \infty}{} 0$$
Or, par hypothèse, pour tout $\code{T}$ de taille $n$, l'appel $\mc{A}(\code{T})$ a une complexité $C_n$ vérifiant~:
$$C_n > \frac{n}{100}$$
à constante multiplicative près. En effet, $C_n$ lit strictement plus d'un centième du tableau, effectuant alors plus de $\frac{n}{100}$ opérations élémentaires. Cela implique que~:
$$\frac{C_n}{n} > \frac{1}{100}$$

Ce qui contredit le fait que $\frac{C_n}{n} \flechelim{n \to + \infty}{} 0$. D'où l'absurdité puis le résultat.

\subsubsection*{Question 2}

Soit $\alpha \in ]0 ; 1[$. On suppose que $C_n \symbolelim{n \to + \infty}{}{\Theta(n^\alpha)}$. Alors il existe $K_1 > 0$ et $K_2 > 0$ tels que~:
$$K_1 n^\alpha \leq C_n \leq K_2 n^\alpha$$
Ainsi, si $n$ est la taille d'un tableau lu à moins de 1\%, alors~:
\begin{align*}
    &C_n \leq \frac{n}{100}\\
    \text{donc} \qquad &K_1 n^\alpha \leq \frac{n}{100}\\
    \text{donc} \qquad &n^{\alpha-1} \leq \frac{1}{100 K_1}\\
    \text{donc} \qquad &\ln(n) \geq \frac{100K_1}{1-\alpha}\\
    \text{donc} \qquad &n \geq \bigg(e^{100K_1}\bigg)^{1-\alpha}
\end{align*}

Par souci de simplification, on peut prendre $K_1 = 1$, ce qui est notamment vérifié quand $C_n \symbolelim{n \to + \infty}{}{\sim} n^\alpha$ (assez commun quand on a un $\Theta$ en étude de complexité).
\begin{itemize}
    \item Pour $\alpha = \frac{1}{2}$, on obtient le minorant $n_\mr{min} = \bigg(e^{100}\bigg)^{1/2} = e^{50}$
    \item Pour $\alpha = 0.9$, on obtient le minorant $n_\mr{min} = \bigg(e^{100}\bigg)^{0.9} = e^{90}$
\end{itemize}

À partir de telles tailles de tableau, l'algorithme ne lira pas plus d'un centième de l'entrée.




\end{document}